\section{Conclusion}
\label{sec:conclusion}

We studied how the order of few-shot demonstrations affects in-context learning across shot counts, task types and model sizes, using one controlled protocol on five \qwen models.
Demonstration order changes performance by about 1--13 points of mean best--worst range, depending strongly on task and model.
Order variance is at its plateau already at two shots; the widely reported rise up to four shots comes from the number of orders sampled.
Position influence is task- and size-dependent rather than a universal recency/primacy split: we see recency for small models on sentiment and on summarization, anti-recency for mid-size models on NLI, and roughly uniform influence for the largest models.
Scale reduces order sensitivity on simple classification but not on \mnli or summarization.
The key takeaway is that whether a ``best order'' exists is itself an empirical question: for strong models on easy tasks it is mostly noise, while on harder tasks it is systematic and worth handling.

Future work should replicate these measurements on other model families and with chat templates; test whether NLI anti-recency comes from contrast with an inconsistent nearest demonstration, by flipping to correct \vsx incorrect labels; use larger test sets and full permutation enumeration for $K\leq5$ to estimate true worst cases; probe content-level primacy in generation; and test whether ensembling over orders recovers the systematic losses on \mnli.
